\chapter{En el exterior}
\label{ch:internationalism}

\section{Por qué una isla pequeña y pobre tenía política exterior}

Durante treinta años Cuba condujo una política exterior de un tamaño varios
órdenes de magnitud superior al suyo, y cualquier relato del período que trate
esto como algo periférico ha pasado por alto cuántos recursos, cuánta
autocomprensión y cuánta reputación del país se fueron en ello. En su punto más
alto Cuba tenía decenas de miles de soldados en África, médicos y maestros en
docenas de países, y una capacidad de atraer la atención del mundo no alineado
que ningún Estado de seis millones de habitantes ha tenido antes ni después.

Tres motivos corrían juntos y no siempre es posible separarlos. Hubo convicción:
la dirigencia creía genuinamente en una revolución global y actuó según esa
creencia a costo real. Hubo seguridad: un país bajo amenaza permanente de una
superpotencia a noventa millas compra aliados y compra estatura, y la estatura es
una forma de disuasión. Y estaba el patrón: después de 1972 los despliegues
militares cubanos en África eran logísticamente imposibles sin transporte y
financiamiento soviéticos, y servían a propósitos soviéticos, lo que no es lo
mismo que decir que Moscú los ordenara ---la intervención en Angola en particular
parece haberse decidido en La Habana y haberse presentado a Moscú después---.

\section{La década guerrillera}

A lo largo de los años sesenta Cuba apoyó la insurgencia armada por toda América
Latina, en Venezuela, Guatemala, Colombia, Perú, Bolivia y Nicaragua, y en
África. Casi todo fracasó. La doctrina, elaborada por Régis Debray a partir de
los escritos de Guevara, sostenía que un pequeño núcleo armado en el campo podía
crear las condiciones de la revolución en vez de esperarlas, lo cual era una
mala lectura del propio caso cubano: la Sierra Maestra había triunfado junto a un
movimiento urbano, un ejército en desintegración y un dictador aislado, nada de
lo cual exigía la doctrina.

Guevara salió de Cuba en 1965 para aplicarla, primero en el Congo, donde la
campaña fue un fracaso que él mismo documentó con notable honestidad, y luego en
Bolivia, donde fue capturado y muerto por el ejército boliviano con asistencia de
la CIA en octubre de 1967. Su muerte lo convirtió en la imagen más reproducida
del siglo XX y retiró de la política cubana a la única figura cuyo prestigio era
independiente del de Castro.

La política le costó también a Cuba con Moscú, que perseguía la coexistencia
pacífica y tenía relaciones de trabajo con varios de los gobiernos que La Habana
intentaba derrocar, y con los partidos comunistas latinoamericanos, a los que
Cuba denunciaba como pasivos. La Conferencia Tricontinental de 1966 en La Habana
fue el punto más alto del tercermundismo cubano y estuvo cerca del punto más bajo
de las relaciones cubano-soviéticas.

La alineación se recompró en 1968. En agosto, cuando los tanques del Pacto de
Varsovia entraron en Checoslovaquia, Castro fue a la televisión y respaldó la
invasión. Lo hizo en un discurso que concedía casi todo ---que la invasión
violaba la legalidad, que el campo socialista no tenía derecho a hacerla bajo las
reglas que él mismo invocaba--- y la respaldaba de todos modos con el argumento
de que la alternativa era la deriva de Checoslovaquia hacia Occidente. El
discurso es el documento más revelador de la política exterior cubana del
período, y el canje que encarnaba fue explícito en cuestión de meses: las entregas
de petróleo soviético aumentaron, siguió la integración económica de los setenta,
y el apoyo cubano a los movimientos guerrilleros en América Latina se fue
retirando.

\section{África}

El compromiso militar de Cuba en África comenzó en Argelia en 1963 y se volvió
grande en Angola en 1975.

Cuando el imperio portugués se derrumbó, la independencia de Angola quedó
disputada entre tres movimientos. El MPLA, alineado con los soviéticos, controlaba
Luanda; el FNLA avanzaba desde Zaire; y la UNITA la respaldaba desde el sur
Sudáfrica, que invadió en octubre de 1975 con una columna blindada rumbo a la
capital. La respuesta cubana, la Operación Carlota \citep{gleijeses2002}, comenzó
en noviembre: tropas transportadas en aviones cubanos y después soviéticos, que
detuvieron el avance sudafricano a corta distancia de Luanda. En meses había
36.000 soldados cubanos en Angola. Se quedaron dieciséis años. A lo largo de todo
el período sirvieron allí unos 300.000 militares y 50.000 civiles cubanos
\citep{gleijeses2013}.\unv{}

La fase decisiva de la guerra llegó en 1987--88 en torno a Cuito Cuanavale, en el
sureste, donde una fuerza cubana y angolana detuvo una ofensiva sudafricana, y
después en un empuje blindado cubano hacia la frontera de Namibia que cambió el
equilibrio militar en el suroeste. El resultado de la batalla se disputa en el
detalle ---ambos bandos la reclaman--- pero su consecuencia no: las negociaciones
cuatripartitas que siguieron produjeron la retirada sudafricana de Angola, la
independencia de Namibia en 1990 y la retirada cubana, y ocurrieron porque
Pretoria concluyó que ya no podía sostener militarmente la línea. Nelson Mandela,
liberado en 1990, fue a La Habana en 1991 y lo dijo directamente. Cualquiera que
sea el resto del balance, la contribución cubana al fin del control sudafricano
sobre Namibia y a la presión que acabó con el apartheid es real, se pagó con
vidas cubanas y no se hizo por dinero.

Etiopía es la otra mitad del registro y es mucho peor. En 1977--78 Cuba envió
unos 16.000 soldados a defender a Etiopía de una invasión somalí del Ogadén
---defendible en sus propios términos, puesto que Somalia era el agresor y Cuba
se negó a sumarse a la ofensiva subsiguiente sobre territorio somalí---. Pero el
gobierno que se defendía era el Derg de Mengistu Haile Mariam, que estaba en
plena Terror Rojo, y la presencia cubana continuó a través de él y a través de la
hambruna y los reasentamientos forzosos de los años ochenta. Cuba había sido
además aliada de Somalia y del movimiento independentista eritreo antes de cambiar
de bando, y los eritreos, que se habían entrenado en Cuba, estaban ahora del otro
lado de los fusiles. Fuera lo que fuese el internacionalismo en Angola, en
Etiopía significó sostener a un cliente asesino en masa del mismo patrón.

El conteo oficial cubano de muertos en todas las misiones internacionalistas,
publicado en 1991, fue de 2.289.\unv{} A las familias se les dijo lo necesario y
los ataúdes volvieron a casa juntos, en una sola ceremonia, dieciséis años de
ellos a la vez, en diciembre de 1989.

\subsection{Los médicos}

Corriendo en paralelo, y sobreviviendo a todo lo demás, estaba el
internacionalismo médico. La primera brigada médica cubana fue a Argelia en 1963,
cincuenta y seis personas, enviadas por un país que acababa de perder la mitad de
sus propios médicos \citep{kirk2015}. Para los años ochenta había médicos cubanos
trabajando en África, Centroamérica y el Caribe, casi siempre sin cobrar, en
lugares a los que ningún otro país mandaba a nadie. La Escuela Latinoamericana de
Medicina, fundada más tarde, formaba gratis a estudiantes extranjeros.

Este es el origen de la reputación que el capítulo~\ref{ch:life} sostendrá que es
el activo intangible más valioso de Cuba ---y, a partir de los años dos mil, del
mayor renglón exportador del país, que es otra cosa y que trae consigo sus
propios problemas de paga, consentimiento y condiciones, de los que la Parte II
se ocupará con honestidad---.

\section{El Mariel}

En abril de 1980 un autobús atravesó la cerca de la embajada del Perú en La
Habana. Los guardias se retiraron tras una disputa sobre el asilo, y en setenta y
dos horas había unas 10.800 personas dentro de los jardines de la embajada. La
escala del asunto ---diez mil personas en un jardín, en un país cuyo gobierno
llevaba veinte años diciendo que solo un puñado de gusanos quería irse--- fue una
sacudida de la que el relato oficial nunca se recuperó del todo.

La respuesta de Castro fue abrir el puerto del Mariel e invitar a los
cubanoamericanos a venir a recoger a sus parientes en barco. Entre abril y
octubre de 1980 cruzaron a Florida unas 125.000 personas. El Estado cubano usó la
apertura para vaciar en las embarcaciones parte de sus prisiones y hospitales
psiquiátricos ---la cifra se disputa, la práctica no, y fue deliberada, tanto
para librarse de una carga como para desacreditar la emigración como criminal---.
Organizó también actos de repudio, manifestaciones de turba frente a las casas de
quienes se iban, en las que se reunía a vecinos y compañeros de trabajo para
gritar, tirar huevos y a veces golpear a quien había anunciado su partida. Se
obligó a cubanos a brutalizar a sus vecinos como condición de su propia posición.

Las consecuencias corrieron en ambas direcciones. En Estados Unidos, el Mariel
llegó en año electoral, desbordó la capacidad de Florida, quedó fijado en la
opinión pública por el contingente criminal, y dañó duraderamente la posición de
los inmigrantes cubanos en general. En Cuba, la memoria de las turbas de repudio
es una de las heridas más agudas que la revolución le infligió a su propia
sociedad ---no por lo que hizo el Estado, sino por lo que hizo que la gente común
se hiciera entre sí, que es más difícil de perdonar y más difícil de olvidar---.
Todo proceso de reconciliación que se proponga en la Parte III tiene que tener una
respuesta para el acto de repudio, y la respuesta no es que fue hace mucho.

\section{El final de la década}

Dos hechos cerraron el período.

En junio y julio de 1989 el general Arnaldo Ochoa, héroe de la campaña de Angola
y uno de los tres o cuatro oficiales más prominentes del país, fue arrestado,
juzgado y ejecutado junto con otros tres, y un grupo de altos funcionarios del
Ministerio del Interior, entre ellos los hermanos de la Guardia, fue condenado por
narcotráfico y corrupción. El Ministerio del Interior fue purgado y puesto bajo
control militar. Que hubo tráfico no se disputa seriamente. Si el juicio trataba
sobre el tráfico, sí: Ochoa era popular, estaba asociado a los oficiales que
habían visto cómo estaba cambiando la Unión Soviética, y según se ha informado
había criticado a la dirigencia. El proceso se televisó, las confesiones tenían la
cadencia de 1971, y jamás ha sido posible una investigación independiente. Lo
cierto es que el asunto retiró a la única figura de Cuba con una reputación
militar comparable a la de los Castro, exactamente en el momento en que se
disolvía el bloque soviético.

En abril de 1989 Mijaíl Gorbachov llegó a La Habana. La visita fue correcta y
fría. Cuba ya había prohibido dos revistas soviéticas, \emph{Novedades de Moscú} y
\emph{Sputnik}, por difundir la perestroika; Castro ya había pronunciado el
discurso que le dio al período su consigna, \emph{socialismo o muerte}. La
posición cubana era que la reforma en la Unión Soviética era asunto soviético y
que Cuba no la seguiría. No la siguió. Diecinueve meses después de la visita de
Gorbachov, el comercio que había sostenido al país durante dieciocho años empezó a
dejar de llegar.

\begin{ledger}{en el exterior, 1963--1990}
\emph{Logrado.} Una contribución militar decisiva a la derrota de la expansión
sudafricana en el África austral, a la independencia de Namibia y a la presión
que terminó con el apartheid ---reconocida por Mandela, pagada con vidas cubanas
y no pagada por nadie---. Brigadas médicas en docenas de países desde 1963,
semilla del activo más duradero de Cuba y de su eventual mayor exportación. Una
estatura en el mundo no alineado desproporcionada al tamaño del país, que
funcionaba como componente real de su seguridad.

\emph{Costo.} Una década de insurgencia exportada que fracasó en todos los
lugares donde se intentó y mató a los que fueron enviados a intentarlo. El
respaldo a la invasión de Checoslovaquia, dado con pleno conocimiento de lo que
era, como precio del petróleo. Un largo despliegue en apoyo del Derg a través del
Terror Rojo y la hambruna. Dieciséis años de guerra librados por reclutas cuyas
familias supieron poco y cuyos muertos volvieron a casa en una sola ceremonia. Y
en casa, en 1980, la movilización deliberada de cubanos comunes para agredir a
sus vecinos que se iban ---lo único de este balance para lo que no cabe alegar
necesidad externa alguna---.
\end{ledger}
